\section*{3.4 Generalization to search problems (author source p.19)}
Let $f$ be a relation, $f\subseteq\bits^n\times V$. For a gadget $g:\bits^m\to\bits$, we define a composition $f\comp g\subseteq(\bits^m)^n\times V$ as a relation, such that
\[(z_1,z_2,\ldots,z_n,r)\in(f\comp g)\iff(g(z_1),g(z_2),\ldots,g(z_n),r)\in f.\]
One can observe that the proofs of Theorem 3.2 and Theorem 3.4 work as well when we let $f$ to be a relation. However, the Corollary 3.5 does not hold for the relations since $\DT(f)$ can be exponentially greater than $\CC(f)$.
\section*{4.1 Conjunction DAGs (author source p.20)}
A conjunction DAG was first defined formally in Garg et al. (2020), though it appears implicitly in previous work. A conjunction DAG over a set of variables is a single-rooted DAG with the following additional information:
\begin{itemize}
\item Each internal vertex is annotated with a variable, and it has two outgoing edges, one labeled 0 and the other labeled 1.
\item Each vertex $v$ is annotated with a partial assignment $\rho(v)$ with the following constraint. Suppose that $v$ queries $x_i$, and the answer $b$ leads to the vertex $v_b$. Then the partial assignment $\rho(v_b)$ is a subset of the partial assignment $\rho(v)\cup\{x_i\leftarrow b\}$. (This is not identical to the definition in Garg et al. (2020), but morally the same.)
\item The partial assignment at the root is the empty assignment.
\item Each leaf is annotated with some value.
\end{itemize}
A decision DAG computes $f$ if for every leaf $\ell$ annotated with $y_\ell$, $\rho(\ell)$ is a $y_\ell$-certificate of $f$.
Every decision tree is a decision DAG. There are three parameters of interest: the (total) size (number of vertices), the leaf size (number of leaves), and the width (maximum number of variables in any $\rho(v)$).
